\documentclass[a4paper,11pt]{article}
\usepackage{exam-style-341}
\examinehead{341 农业综合知识三 · 数据库分章练习}{第05章\quad 事务、并发与恢复}
\begin{document}

\unitcover{341 农业综合知识三 · 数据库分章练习}{第05章\quad 事务、并发与恢复}{练习（共 100 分）}{约 60 分钟}{事务的概念与 ACID 性质；并发操作带来的丢失修改、不可重复读和读脏数据；封锁与共享锁、排他锁；三级封锁协议；两段锁协议与可串行化调度；死锁与活锁；故障类型与恢复技术（日志、检查点、UNDO/REDO）；数据库安全性（GRANT、REVOKE、视图）。}

\begin{knowmap}
\textbf{事务与 ACID：}事务是用户定义的一个数据库操作序列，要么全做、要么全不做；具有原子性、一致性、隔离性和持久性\\[2pt]
\textbf{并发三类问题：}丢失修改（两个事务读同一数据并修改，后提交者覆盖前者）· 不可重复读（两次读同一数据结果不同）· 读脏数据（读到未提交且被回滚的数据）\\[2pt]
\textbf{封锁：}排他锁 X 锁（加锁后可读写，其他事务不能加任何锁）· 共享锁 S 锁（加锁后只能读，其他事务可再加 S 锁）；三级封锁协议逐级加强并发控制\\[2pt]
\textbf{两段锁协议：}扩展阶段只能加锁不能解锁，收缩阶段只能解锁不能加锁；遵守两段锁协议的并发调度一定是可串行化的\\[2pt]
\textbf{故障与恢复：}事务故障（UNDO）· 系统故障（UNDO 未提交事务、REDO 已提交事务）· 介质故障（重装数据库副本 + 重做日志）；日志记录事务对数据的修改，检查点减少恢复时需要重做的事务\\[2pt]
\textbf{安全性：}自主存取控制通过 GRANT 和 REVOKE 授权与收回权限；视图可以把用户限定在其权限内的数据上
\end{knowmap}

\begin{keyterms}
\textbf{事务}：用户定义的一个数据库操作序列，这些操作要么全做、要么全不做，是一个不可分割的工作单位。\\
\textbf{原子性}：事务中的操作要么全部执行，要么全部不执行；若中途失败，已做的修改要全部撤销。\\
\textbf{一致性}：事务执行的结果必须使数据库从一个一致性状态变到另一个一致性状态。\\
\textbf{隔离性}：一个事务的执行不能被其他事务干扰，并发执行的事务之间不能互相看到未提交的中间结果。\\
\textbf{持久性}：事务一旦提交，它对数据库中数据的改变就是永久性的，后续操作或故障不应影响其结果。\\
\textbf{丢失修改}：两个事务读入同一数据并加以修改，后提交的事务覆盖了先提交事务的修改，导致先前的修改丢失。\\
\textbf{不可重复读}：事务先后两次读取同一数据，由于其他事务的修改、插入或删除，两次结果不一致。\\
\textbf{读脏数据}：事务读到另一个事务修改但尚未提交的数据，而该事务随后被回滚，所读数据是错误数据。\\
\textbf{排他锁}：又称 X 锁或写锁，事务对数据加 X 锁后可以读写该数据，其他事务不能对该数据加任何锁。\\
\textbf{共享锁}：又称 S 锁或读锁，事务对数据加 S 锁后只能读不能写，其他事务可以再对该数据加 S 锁但不能加 X 锁。\\
\textbf{三级封锁协议}：在运用 X 锁和 S 锁时约定不同加锁规则的三级协议；一级防止丢失修改，二级进一步防止读脏数据，三级进一步防止不可重复读。\\
\textbf{两段锁协议}：所有事务必须分两个阶段对数据项加锁和解锁：扩展阶段只能加锁不能解锁，收缩阶段只能解锁不能加锁。\\
\textbf{死锁}：两个或多个事务都已封锁了一些数据对象，又都请求对方已封锁的对象，导致这些事务永远等待、无法继续执行。\\
\textbf{可串行化调度}：多个事务的并发执行结果与按某种顺序串行执行这些事务的结果相同，则该并发调度是可串行化的，这是并发事务正确性的准则。
\end{keyterms}

\begin{contrast}{原子性 / 一致性 / 隔离性 / 持久性（每行依次为：含义·主要保证手段·被破坏时的典型现象）}
\citem{原子性}{事务中的操作要么全做、要么全不做}{日志与 UNDO 机制}{事务中途失败，部分修改留在数据库中}
\citem{一致性}{事务使数据库从一个一致性状态变到另一个一致性状态}{完整性约束检查、应用逻辑}{转账只扣款不入账，总额不守恒}
\citem{隔离性}{并发事务之间互不干扰}{封锁、两段锁协议、可串行化调度}{读到其他事务未提交的数据或中间结果}
\citem{持久性}{事务提交后的修改永久有效}{日志与 REDO 机制、检查点}{提交后系统故障导致修改丢失}
\end{contrast}

\begin{contrast}{丢失修改 / 不可重复读 / 读脏数据（每行依次为：现象·产生原因·主要解决手段）}
\citem{丢失修改}{两个事务修改同一数据，后提交者覆盖前者}{两个事务读入同一数据后各自修改并写回}{一级封锁协议：修改前必须先加 X 锁并保持到事务结束}
\citem{不可重复读}{同一事务两次读取同一数据结果不同}{其他事务在两次读之间修改、插入或删除了数据}{三级封锁协议：读数据时加 S 锁并保持到事务结束}
\citem{读脏数据}{读到其他事务未提交且后来被回滚的数据}{读操作未加锁，读到了未提交的中间值}{二级封锁协议：读数据前加 S 锁，读完即释放}
\end{contrast}

\begin{contrast}{共享锁 / 排他锁（每行依次为：允许的操作·与其他锁的相容性·加锁规则）}
\citem{共享锁}{加锁事务只能读，不能写}{与其他 S 锁相容，与 X 锁不相容}{读取数据前申请 S 锁}
\citem{排他锁}{加锁事务可以读，也可以写}{与任何锁都不相容}{修改数据前申请 X 锁，并通常保持到事务结束}
\end{contrast}

\begin{contrast}{事务故障 / 系统故障 / 介质故障（每行依次为：产生原因·影响范围·恢复方法）}
\citem{事务故障}{运算溢出、死锁被选中撤销等使事务未能正常结束}{只影响该事务}{反向扫描日志，对未提交事务执行 UNDO}
\citem{系统故障}{断电、操作系统崩溃、DBMS 异常终止等}{内存中数据丢失，磁盘数据未损坏}{UNDO 未提交事务，REDO 已提交事务}
\citem{介质故障}{磁盘损坏、磁头碰撞、强磁场干扰等}{磁盘上的数据部分或全部损坏}{重装最近数据库副本，再用日志 REDO 已提交事务}
\end{contrast}

\qtypename{一、单项选择题}{10 题，每题 2 分，共 20 分}
\q 事务的 ACID 性质中，“事务一旦提交，它对数据库中数据的改变就是永久性的”指的是（\quad）
\choicesline{A. 原子性}{B. 一致性}{C. 隔离性}{D. 持久性}

\q 事务 $T_1$ 修改了数据 $A$ 但尚未提交，事务 $T_2$ 读取了 $A$ 的值，随后 $T_1$ 被回滚，$T_2$ 读到的数据是错误的。这种现象称为（\quad）
\choicesline{A. 丢失修改}{B. 不可重复读}{C. 读脏数据}{D. 死锁}

\q 事务对数据加了排他锁（X 锁）后，其他事务（\quad）
\choices{可以对该数据加 S 锁}{可以对该数据加 X 锁}{不能对该数据加任何锁}{只能读不能写}

\q 一级封锁协议能够防止（\quad）
\choicesline{A. 丢失修改}{B. 读脏数据}{C. 不可重复读}{D. 死锁}

\q 两段锁协议规定事务（\quad）
\choices{先加锁后解锁，且加锁阶段不能解锁}{只能加一次锁}{只能加 S 锁}{加锁和解锁可以任意交替}

\q 下列措施中，属于死锁预防方法的是（\quad）
\choices{一次封锁法或顺序封锁法}{超时法}{事务等待图法}{检查点法}

\q 数据库系统发生事务故障时，恢复子系统对未提交事务的处理是（\quad）
\choicesline{A. 执行 UNDO}{B. 执行 REDO}{C. 重装数据库副本}{D. 不做任何处理}

\q 系统故障恢复时，对故障发生前已提交的事务应（\quad）
\choicesline{A. 执行 UNDO}{B. 执行 REDO}{C. 删除其日志}{D. 重新执行该事务的用户程序}

\q 日志文件中设置检查点的主要目的是（\quad）
\choices{减少恢复时需要扫描的日志范围和需要重做的事务}{提高查询速度}{防止死锁}{保证事务的隔离性}

\q 语句 GRANT SELECT ON Student TO U1 属于数据库的（\quad）
\choicesline{A. 安全性控制}{B. 完整性控制}{C. 并发控制}{D. 恢复控制}

\qtypename{二、多项选择题}{5 题，每题 4 分，共 20 分；每题有两个或两个以上正确项}
\q 事务的 ACID 性质包括（\quad）
\choicesline{A. 原子性}{B. 一致性}{C. 隔离性}{D. 持久性}

\q 并发操作可能带来的数据不一致性问题有（\quad）
\choicesline{A. 丢失修改}{B. 不可重复读}{C. 读脏数据}{D. 数据冗余}

\q 关于封锁，下列说法正确的有（\quad）
\choices{共享锁与共享锁相容}{共享锁与排他锁相容}{排他锁与排他锁不相容}{排他锁与共享锁不相容}

\q 数据库系统可能发生的故障类型有（\quad）
\choicesline{A. 事务故障}{B. 系统故障}{C. 介质故障}{D. 计算机病毒故障}

\q 关于日志文件，下列说法正确的有（\quad）
\choices{日志记录事务对数据库的每一次更新操作}{日志记录中应包含事务开始、提交和回滚标记}{利用日志可以执行 UNDO 和 REDO}{日志必须写在数据库之前（先写日志）}

\qtypename{三、判断题}{5 题，每题 2 分，共 10 分；正确的画“√”，错误的画“×”并写出改正后的表述}
\q 事务中的操作要么全做、要么全不做，体现原子性。\quad（\quad）\quad 改正：\fixblank

\q 遵守两段锁协议的并发调度一定是可串行化的。\quad（\quad）\quad 改正：\fixblank

\q 死锁属于介质故障。\quad（\quad）\quad 改正：\fixblank

\q 系统故障恢复时，对已提交事务应执行 UNDO。\quad（\quad）\quad 改正：\fixblank

\q 读脏数据是指事务读到了另一个事务已提交的数据。\quad（\quad）\quad 改正：\fixblank

\qtypename{四、填空题}{5 题，每题 2 分，共 10 分}
\q 事务的 ACID 性质是指原子性、\fillblank{}、隔离性和 \fillblank{}。

\q 并发操作带来的三类数据不一致性问题是丢失修改、\fillblank{} 和 \fillblank{}。

\q 封锁的基本类型有 \fillblank{} 锁和 \fillblank{} 锁，其中前者又称写锁，后者又称读锁。

\q 三级封锁协议中，一级封锁协议可以防止 \fillblank{}，二级封锁协议可以进一步防止 \fillblank{}。

\q 数据库系统通常将故障分为事务故障、\fillblank{} 和 \fillblank{} 三类。

\qtypename{五、简答题}{2 题，每题 8 分，共 16 分}
\q 简述事务的 ACID 性质，并说明并发操作可能带来的三类数据不一致性问题。\anslines{2.2cm}

\q 简述三级封锁协议的内容，并分别说明每一级封锁协议能够解决什么问题。\anslines{2.2cm}

\qtypename{六、综合题}{2 题，每题 12 分，共 24 分}
\q 某银行数据库中有账户表 Account(Ano, Balance)，账户 $A$ 的余额为 1000 元，现执行从账户 $A$ 向账户 $B$ 转账 100 元的业务。请完成：
\begin{enumerate}[label=（\arabic*）, leftmargin=2.4em, itemsep=1pt]
  \item 用 SQL 事务控制语句写出该转账业务的事务框架，要求体现事务的原子性。（3 分）
  \item 判断下列三个并发执行场景分别会产生哪一类数据不一致性问题，并简要说明理由：（6 分）
  \begin{itemize}[leftmargin=1.6em, itemsep=1pt]
    \item 场景一：$T_1$ 读 $A=1000$，$T_2$ 读 $A=1000$，$T_1$ 把 $A$ 改为 900 并写回，$T_2$ 把 $A$ 改为 800 并写回。
    \item 场景二：$T_1$ 读 $A=1000$，$T_2$ 把 $A$ 改为 900 并提交，$T_1$ 再次读 $A$ 得 900。
    \item 场景三：$T_1$ 把 $A$ 改为 900 但尚未提交，$T_2$ 读到 $A=900$，随后 $T_1$ 回滚。
  \end{itemize}
  \item 说明两段锁协议的内容，并指出遵守两段锁协议的并发调度具有什么性质。（3 分）
\end{enumerate}
\anslines{4.5cm}

\q 某数据库系统在运行过程中发生故障，故障发生时的日志记录（按时间先后）如下（$A,B,C,D$ 为数据项，括号内为修改前的值和修改后的值）：
\begin{center}
\small
\begin{tabular}{@{}ll@{}}
\toprule
日志记录 & 说明 \\
\midrule
$<$T1, START$>$ & 事务 T1 开始 \\
$<$T1, A, 100, 200$>$ & T1 把 A 由 100 改为 200 \\
$<$T1, COMMIT$>$ & 事务 T1 提交 \\
$<$T2, START$>$ & 事务 T2 开始 \\
$<$T2, B, 50, 80$>$ & T2 把 B 由 50 改为 80 \\
$<$CKPT (T1)$>$ & 建立检查点，此时 T1 已提交，T2 正在执行 \\
$<$T3, START$>$ & 事务 T3 开始 \\
$<$T3, C, 30, 60$>$ & T3 把 C 由 30 改为 60 \\
$<$T2, COMMIT$>$ & 事务 T2 提交 \\
$<$T4, START$>$ & 事务 T4 开始 \\
$<$T4, D, 10, 20$>$ & T4 把 D 由 10 改为 20 \\
$<$T4, COMMIT$>$ & 事务 T4 提交 \\
故障发生 & 系统崩溃，内存中的信息丢失 \\
\bottomrule
\end{tabular}
\end{center}

请完成：
\begin{enumerate}[label=（\arabic*）, leftmargin=2.4em, itemsep=1pt]
  \item 指出故障发生时 T1、T2、T3、T4 各自的事务状态。（3 分）
  \item 写出恢复时需要执行 UNDO 和 REDO 的事务集合。（3 分）
  \item 写出恢复后数据项 $A,B,C,D$ 的值。（3 分）
  \item 说明检查点在恢复过程中的作用。（3 分）
\end{enumerate}
\anslines{4.5cm}

\end{document}
